\subsection{Representation}

Phase I demonstrates that a minimally specified spiking substrate
develops sensory-specific organization (afferent selectivity
\(\pm 0.8\)--\(0.9\); 91\% run-identity decode of per-neuron afferent
vectors; repeatable per-presentation count vectors). It equally
demonstrates that the transient representation is weak where it
matters: cross-pattern response separation stalled near 0.67--0.74
against a 0.60 bar, and the two dynamics interventions (adaptation,
lateral inhibition) moved it in the wrong direction
\cite{e1proto,e3bproto}. The binding problem of this organism is
representational: 24 afferents drive one 40-neuron recurrent pool
through random sparse wiring, so every readout mixes patterns. Nothing
in the record suggests this is fixable by further rate dynamics; the
only lever that has ever moved coexistence is budget partition (V2.3,
blocked-only) \cite{overview,synmem}.

\subsection{Persistence}

Persistence has two measured faces. Structural persistence is real and
bounded by mechanics: weights survive silence (matrix drift \(\sim
1.4\%\)/20 s), decay is the dominant eraser (SDE-C2), and bounded
protection stops the runaway that unbounded protection caused
\cite{bbounded,cllarecord}. Activity persistence is information-poor:
the V2.1 slow state self-sustains but saturates, and its endogenous
form is deterministic pacemaking \cite{infoaudit}. The paper's
strongest persistence claim is negative and precisely delimited: no
persistent substrate measured in Phase I retains episode identity
through repeated or alternated exposure --- including, in the testable
case, during the stimulus itself.

\subsection{Interference}

Interference is order-dependent and budget-mediated. Isolated
training selects; blocked training overwrites; alternation superposes
into a third configuration. The causal probes attribute the erasure to
passive decay acting through M2's shrinking working budget, with M4 as
the executing pruner; STDP's \(a^-/a^+\) asymmetry is secondary
(SDE-D). Protection fixes stability, not interference: even with
bounded protected structure, blocked order captures the budget for the
latest block \cite{capability}. Interference is therefore best
described as an \emph{allocation} phenomenon --- competition for a
shared local budget --- rather than a plasticity-sign phenomeno.

\subsection{Allocation}

The allocation sequence is the paper's procedural core: five
registered experiments, each gated byte-identical, each moving the
identified bottleneck one step downstream (gate \(\to\) trigger
semantics \(\to\) visibility \(\to\) source \(\to\) post-permanence
growth). The final measured residual is that the working substrate of a
not-yet-presented pattern is dismantled before its first exposure, and
that even a corrected allocator cannot grow the protected structure
once substrate is gone. This is the point where Phase I's causal
narrowing arrives at the architectural frontier: allocation works when
the pattern is active, and fails for patterns that have never been
active --- which is exactly the class that sequential memory requires.

\subsection{Recruitment}

The recruitment deficit is the frontier's content: a genuinely new
late-arriving pattern cannot recruit postsynaptic activity because its
afferent substrate was churned away, and no membrane-side,
input-proportional gain can amplify a substrate that is not there. The
IL control proves the stimulus is not at fault. Two registered
mechanism tests bounded the hypothesis without resolving it; the
threshold-completion cap fixed the aggregate operating point but could
not repair first-exposure rate or substrate survival. The discussion
stops here by intent: Phase I's evidence supports a direction
(architectures for multiple independently retrievable persistent
representations under sequential experience), not a design.